%%%%%%%%%%%%%%%%%%%%%%%%%%%%%%%%%%%%%%%%%%%%%%%%%%%%%%%%%%%%%%%%%%%%%%
\section{The Static Metric from Information Propagation}
\label{sec:metric}
%%%%%%%%%%%%%%%%%%%%%%%%%%%%%%%%%%%%%%%%%%%%%%%%%%%%%%%%%%%%%%%%%%%%%%

We now derive the full spacetime metric in the static, spherically symmetric limit.
The previous section established the static solution $q(r)=e^{-GM/rc^2}$ from the field equation $\nabla^2(\ln 1/q)=0$.
This section shows that the same information-theoretic principles that fix the temporal metric component also fix the spatial components, yielding a complete metric with precise post-Newtonian predictions.

\subsection{Causal Locking and the Temporal Metric}

The global projection clock advances at rate $1/t_P$ in coordinate time $t$.
A cell $i$ with free-capacity fraction $q_i$ can accept new information only when it is free.
The fraction of coordinate time during which the cell is receptive to causal updates is exactly $q_i$.
Therefore, the proper time $d\tau$ experienced at a point with free capacity $q$ is related to coordinate time $dt$ by
\begin{equation}
\boxed{d\tau = q\,dt}.
\label{eq:causal_locking}
\end{equation}
We term this \textbf{causal locking}: the local clock rate equals the free capacity fraction.

In the rest frame of a static observer, the proper time interval is $d\tau^2 = -g_{00}\,dt^2$.
Equation~\eqref{eq:causal_locking} therefore fixes the temporal metric component:
\begin{equation}
\boxed{g_{00} = -q^2}.
\label{eq:g00}
\end{equation}

For $q\to1$ (empty space, far from all masses), $g_{00}\to-1$, recovering the Minkowski form.
For $q<1$ (massive region), $|g_{00}|<1$ and clocks run slower, as required.

\subsection{Information Propagation Constraint and the Spatial Metric}

The spatial metric follows from an equally fundamental information-theoretic constraint.
Consider how information propagates on the cell graph.

Each cell-to-cell hop of information occurs at the fundamental speed $c=a/t_P$, where $a$ is the cell spacing and $t_P=a/c$ is the Planck time.
However, a hop from cell $i$ to cell $j$ can succeed \emph{only if cell $j$ is free}---if cell $j$ is occupied (its 1-bit capacity is saturated), the information must wait for the cell to become available.

In the steady state, cell $j$ is free with probability $q_j$.
The number of attempts needed for a successful hop follows a geometric distribution with mean $1/q_j$.
Therefore the mean hop time is
\begin{equation}
\Delta t_{\rm hop} = \frac{t_P}{q}.
\label{eq:hop_time}
\end{equation}
(For a smooth $q$ field, $q_i\approx q_j\approx q$, and the harmonic mean is well approximated by the local $q$.)

The effective propagation speed across the cell graph is therefore
\begin{equation}
c_{\rm eff} = \frac{a}{\Delta t_{\rm hop}} = c\,q.
\label{eq:ceff}
\end{equation}
In regions with low $q$ (near massive bodies, where most cells are occupied), information propagates more slowly because it must queue behind occupied cells.

The proper distance $dl_{\rm proper}$ traversed by an information signal in proper time $d\tau$ is $dl_{\rm proper}=c\,d\tau_{\rm prop}$, where $d\tau_{\rm prop}$ is the proper time for the signal to cross coordinate distance $dl_{\rm coord}$.
Since $c_{\rm eff}=dl_{\rm coord}/d\tau_{\rm prop}$, we have $d\tau_{\rm prop}=dl_{\rm coord}/(c q)$, and thus
\begin{equation}
\boxed{dl_{\rm proper} = \frac{dl_{\rm coord}}{q}}.
\label{eq:proper_length}
\end{equation}

This is the central result: the effective proper distance is the coordinate distance \emph{stretched} by factor $1/q$, because each cell-to-cell hop requires, on average, $1/q$ attempts.
Equation~\eqref{eq:proper_length} is the spatial analog of Eq.~\eqref{eq:causal_locking}: temporal intervals contract by $q$, spatial intervals dilate by $1/q$---a natural duality rooted in the same information bottleneck.

In isotropic coordinates, where the spatial part of the metric is conformally flat,
\begin{equation}
dl_{\rm coord}^2 = dr^2 + r^2(d\theta^2 + \sin^2\theta\,d\varphi^2),
\end{equation}
the spatial metric components are
\begin{equation}
\boxed{g_{rr} = \frac{1}{q^2},\qquad
g_{\theta\theta} = \frac{r^2}{q^2},\qquad
g_{\varphi\varphi} = \frac{r^2\sin^2\theta}{q^2}}.
\label{eq:spatial_metric}
\end{equation}

\subsection{Full Static Metric}

Combining Eqs.~\eqref{eq:g00} and~\eqref{eq:spatial_metric}, the full line element in isotropic coordinates is
\begin{equation}
\boxed{ds^2 = -q^2\,dt^2 + \frac{1}{q^2}\bigl(dr^2 + r^2 d\Omega^2\bigr)}.
\label{eq:full_metric}
\end{equation}
The metric is fully determined by the single scalar field $q(\mathbf{x})$.

Inserting the static point-mass solution $q(r)=e^{-GM/rc^2}$ established in Section~\ref{sec:static}:
\begin{equation}
\boxed{ds^2 = -e^{-2GM/rc^2}\,dt^2 + e^{2GM/rc^2}\bigl(dr^2 + r^2 d\Omega^2\bigr)}.
\label{eq:DGF_metric}
\end{equation}

Several properties are immediately apparent:
\begin{enumerate}
    \item As $r\to\infty$, $q\to1$, and the metric approaches Minkowski: $ds^2\to -dt^2+dr^2+r^2d\Omega^2$.
    \item Both $g_{00}$ and $g_{rr}$ depend only on $q(r)$, so the metric is determined by a single degree of freedom---no independent temporal and spatial potentials.
    \item The metric is everywhere non-singular for $r>0$: $g_{00}=-e^{-2GM/rc^2}<0$ and $g_{rr}=e^{2GM/rc^2}>0$.
    \item At $r=0$, $g_{00}\to0$ and $g_{rr}\to\infty$, indicating the Planck-scale limit where $q\to0$ and the capacity bound is saturated.
\end{enumerate}

\subsection{Post-Newtonian Expansion}

Define the Newtonian potential in dimensionless form:
\begin{equation}
U(r) \equiv \frac{GM}{rc^2}.
\label{eq:Udef}
\end{equation}
For Solar System applications, $U\lesssim 10^{-5}$ everywhere outside the Sun, so a post-Newtonian expansion in powers of $U$ is well-justified.

Expand the DGF metric~\eqref{eq:DGF_metric}:
\begin{align}
g_{00} &= -e^{-2U} = -\Bigl(1 - 2U + 2U^2 - \frac{4}{3}U^3 + \frac{2}{3}U^4 - \cdots\Bigr), \label{eq:g00_expand} \\
g_{rr} &= e^{2U} = 1 + 2U + 2U^2 + \frac{4}{3}U^3 + \frac{2}{3}U^4 + \cdots. \label{eq:grr_expand}
\end{align}

In the standard parameterized post-Newtonian (PPN) formalism~\cite{Will:1993,Will:2014}, the isotropic metric is written as
\begin{align}
g_{00} &= -1 + 2U - 2\beta U^2 + \mathcal{O}(U^3), \label{eq:ppn_g00} \\
g_{ij} &= (1 + 2\gamma U)\delta_{ij} + \mathcal{O}(U^2). \label{eq:ppn_gij}
\end{align}

Comparing Eq.~\eqref{eq:g00_expand} with Eq.~\eqref{eq:ppn_g00}:
\begin{equation}
\boxed{\beta = 1}.
\label{eq:beta}
\end{equation}
Comparing Eq.~\eqref{eq:grr_expand} with Eq.~\eqref{eq:ppn_gij}:
\begin{equation}
\boxed{\gamma = 1}.
\label{eq:gamma}
\end{equation}

Both PPN parameters take their general-relativistic values.
DGF therefore reproduces GR at first post-Newtonian order.

\subsection{Consistency with Solar System Tests}

With $\beta=\gamma=1$, all three classical tests of general relativity are satisfied identically:

\textbf{Perihelion precession.}
The PPN expression for the perihelion advance per orbit is~\cite{Will:1993}
\begin{equation}
\Delta\phi = \frac{2+2\gamma-\beta}{3}\,\Delta\phi_{\rm GR},
\end{equation}
where $\Delta\phi_{\rm GR}=6\pi GM/[ac^2(1-e^2)]$ is the standard GR result in terms of the semi-major axis $a$ and eccentricity $e$.
With $\beta=\gamma=1$, the prefactor is $(2+2-1)/3=1$, and DGF gives exactly the GR prediction:
\begin{equation}
\boxed{\Delta\phi_{\rm DGF} = \Delta\phi_{\rm GR} = 42.98''\text{/century (Mercury)}}.
\end{equation}

\textbf{Light deflection.}
The PPN expression for the deflection angle of light passing a mass $M$ with impact parameter $b$ is
\begin{equation}
\Delta\theta = \frac{1+\gamma}{2}\,\frac{4GM}{bc^2}.
\end{equation}
With $\gamma=1$, this gives $\Delta\theta = 4GM/(bc^2)$, matching the GR result of $1.75''$ for starlight grazing the solar limb.

\textbf{Shapiro time delay.}
The PPN expression for the round-trip radar delay is
\begin{equation}
\Delta t = \frac{1+\gamma}{2}\,\Delta t_{\rm GR},
\end{equation}
which again reduces to the GR value for $\gamma=1$.

Thus DGF passes all Solar System tests at the current level of precision.
The previous concern that DGF might be ruled out by perihelion data~\cite{...} was based on an incorrect acceleration law $\mathbf{a}=-c^2\nabla q$; the correct metric derivation via causal locking and the propagation constraint yields $\beta=\gamma=1$ and full consistency with observation.

\subsection{Second-Order Deviation from General Relativity}

Although DGF and GR agree at first PPN order, they differ at second post-Newtonian order.
For GR in isotropic coordinates, the exact Schwarzschild metric is~\cite{MTW:1973}
\begin{equation}
ds^2_{\rm GR} = -\left(\frac{1-GM/2rc^2}{1+GM/2rc^2}\right)^2 dt^2 + \left(1+\frac{GM}{2rc^2}\right)^4 (dr^2 + r^2d\Omega^2).
\end{equation}
Expanding:
\begin{align}
g_{00}^{\rm GR} &= -1 + 2U - 2U^2 + \frac{3}{2}U^3 + \mathcal{O}(U^4), \label{eq:gr_g00} \\
g_{rr}^{\rm GR} &= 1 + 2U + \frac{3}{2}U^2 + \frac{1}{2}U^3 + \mathcal{O}(U^4). \label{eq:gr_grr}
\end{align}

Comparing with the DGF expansions~\eqref{eq:g00_expand}--\eqref{eq:grr_expand}:

\begin{center}
\begin{tabular}{c|cccc}
\toprule
Component & Theory & $\mathcal{O}(U)$ & $\mathcal{O}(U^2)$ & $\mathcal{O}(U^3)$ \\
\midrule
\multirow{2}{*}{$-g_{00}$} & DGF & $1-2U$ & $+2U^2$ & $-\frac{4}{3}U^3$ \\
                            & GR  & $1-2U$ & $+2U^2$ & $-\frac{3}{2}U^3$ \\
\midrule
\multirow{2}{*}{$g_{rr}$}   & DGF & $1+2U$ & $+2U^2$ & $+\frac{4}{3}U^3$ \\
                            & GR  & $1+2U$ & $+\frac{3}{2}U^2$ & $+\frac{1}{2}U^3$ \\
\bottomrule
\end{tabular}
\end{center}

The first deviation appears at $\mathcal{O}(U^2)$ in $g_{rr}$:
\begin{equation}
\Delta g_{rr} \equiv g_{rr}^{\rm DGF} - g_{rr}^{\rm GR}
= \frac{1}{2}U^2 + \mathcal{O}(U^3).
\label{eq:delta_grr}
\end{equation}

For the Sun at its surface ($U_\odot \approx 2.1\times10^{-6}$):
\begin{equation}
\Delta g_{rr}(R_\odot) \approx 2.2\times10^{-12}.
\end{equation}

For Mercury's orbit ($U_{\rm Mercury}\approx 2.5\times10^{-8}$):
\begin{equation}
\Delta g_{rr}(r_{\rm Mercury}) \approx 3.1\times10^{-16}.
\end{equation}

The perihelion precession deviation at second PPN order is suppressed by an additional factor of $U$:
\begin{equation}
\frac{\delta(\Delta\phi)}{\Delta\phi} \sim \mathcal{O}(U^2) \sim 6\times10^{-16}\quad\text{(Mercury)}.
\end{equation}
This corresponds to a deviation of $\sim 3\times10^{-14}$ arcseconds per century, compared to the current measurement precision of $\sim 10^{-3}$ arcseconds per century from MESSENGER data~\cite{Park:2017}.
The second-order DGF deviation is therefore \textbf{eight orders of magnitude} below current observational sensitivity.

\subsection{Sensitivity Required for Detection}

To discriminate DGF from GR at the second PPN order would require:

\begin{itemize}
    \item \textbf{Solar spectroscopy:} measuring the gravitational redshift at the solar surface to $\sim 10^{-12}$ fractional precision. Current best: $\sim 10^{-6}$ (GRAVITY/VLTI, PHI/Solar~Orbiter). Required improvement: ${\sim}10^{6}\times$.

    \item \textbf{Planetary ranging:} measuring Mercury's orbit to $\sim 10^{-16}$ fractional precision in the gravitational potential. Current best: $\sim 10^{-8}$ (MESSENGER, BepiColombo). Required improvement: ${\sim}10^{8}\times$.

    \item \textbf{Pulsar timing:} the double pulsar PSR~J0737--3039 has $U\sim 2\times10^{-6}$ at periastron, giving $\Delta g_{rr}\sim 2\times10^{-12}$. The periastron advance is measured to $\sim 10^{-6}$ fractional precision~\cite{Kramer:2006}. Required improvement: ${\sim}10^{6}\times$.

    \item \textbf{Lunar laser ranging:} $U_{\rm Earth}\sim 7\times10^{-10}$ at the Moon's orbit. $\Delta g_{rr}\sim 2.5\times10^{-19}$, requiring ranging precision of $\sim 10^{-19}$---beyond any foreseeable technology.
\end{itemize}

\textbf{Strong-field regime.}
For a neutron star surface where $U\sim 0.1$--$0.2$, the second-order effect reaches $\Delta g_{rr}/g_{rr}\sim 0.5\%$--$2\%$.
However, in this regime the PPN expansion itself breaks down, and one must solve the full nonlinear DGF field equations.
The soft boundary replacing the classical event horizon (Section~\ref{sec:predictions}) provides a qualitatively distinct, strong-field discriminant that does not require precision metrology.

\textbf{Cosmological regime.}
If $q_\infty$ evolves cosmologically (e.g., $q_\infty(t)\propto 1/a(t)$), the cumulative effect over $\sim 10^{10}$ years could produce $\mathcal{O}(1)$ differences in $g_{rr}$ at cosmological distances, potentially observable through integrated Sachs-Wolfe effect anomalies, distance-ladder tensions, or $f\sigma_8$ growth-rate measurements.
This remains an open direction.

\subsection{Comparison of Derivation Methods}

It is instructive to contrast the present derivation with the earlier DGF effective acceleration interface~\cite{DGF_v3_1}.
In that approach, the complex quadratic representation $Q_q=e^{i\pi(1-q)}$ was projected to its real part, and an effective Christoffel symbol was computed to give $\mathbf{a}=-(c^2\pi^2/2)q\nabla q$.
This produced a non-standard PPN $\gamma$ and led to the false conclusion that DGF was ruled out by perihelion data.

The present derivation avoids this detour entirely:
\begin{enumerate}
    \item Causal locking gives $g_{00}=-q^2$ directly, without invoking complex representations.
    \item The information propagation constraint gives $g_{rr}=1/q^2$ directly, without introducing Christoffel symbols or geodesic equations.
    \item Together, the metric is fully determined by $q(r)$ alone, and the PPN expansion follows mechanically.
    \item The result $\beta=\gamma=1$ is \emph{exact} at first PPN order for \emph{any} spherically symmetric $q(r)$, provided the metric takes the form~\eqref{eq:full_metric}.
\end{enumerate}

The key insight is that \textbf{the metric is not a separate postulate}---it is a direct consequence of how information propagates on the cell graph.
Causal locking and the information bottleneck are two faces of the same capacity bound.

\subsection{Connection to the Planck-Cell Paradigm}

The present framework shares a foundational assumption with the Quantum Memory Matrix (QMM) hypothesis of Neukart et al.~\cite{Neukart:2024}:
spacetime at the Planck scale consists of discrete cells, each carrying a finite-dimensional Hilbert space, and physical processes are fundamentally information-encoding operations on this substrate.
In QMM, each Planck cell acts as a quantum memory unit that stores and releases information via ``quantum imprints,'' preserving unitarity in black hole evaporation.
In DGF, each Planck cell acts as an information capacity unit whose occupation state determines the local metric structure.

The independent convergence of these two programs on the Planck-cell paradigm---from opposite directions (black hole information paradox vs.\ decoherence-to-gravity emergence)---lends mutual credibility.
QMM verifies that Planck-scale information storage is mathematically consistent with known physics; DGF provides the dynamical bridge from that information substrate to the classical gravitational field.

\subsection{Summary}

The static metric derivation closes a structural gap in the DGF/LP31 framework.
The chain from axioms to observable metric is now complete:
\begin{equation}
\text{A1+A2}\;\to\;q(r)=e^{-GM/rc^2}\;\to\;g_{00}=-q^2,\;g_{rr}=1/q^2\;\to\;\beta=\gamma=1.
\end{equation}
DGF reproduces all three classical Solar System tests at the precision with which they have been measured.
The second-order deviation from GR, $\Delta g_{rr}\simeq\frac{1}{2}(GM/rc^2)^2$, is a genuine prediction that distinguishes the two theories, albeit at a level (${\sim}10^{-16}$ for Mercury) far below current observational reach.
Qualitative discriminants in the strong-field regime (soft boundary, modified gravitational wave ringdown) offer more promising near-term tests.
\end{document}
